\section{Background}
\label{sec:background}

vFree is composed of two primary strategies which, when combined, make it possible to achieve near-native performance for network-heavy workloads while still maintaining the security benefits of memory protection provided by the vIOMMU. This is done by rewriting the invalidation routines to be more efficient (and configurable), improving support for contiguous page frees, and allowing for deferred free page (DFP) which allows any drivers that use the feature to register a callback that is invoked when a page is unmapped. This allows for the driver to work on other tasks while the unmapping is being processed, which greatly reduces locking overhead and still allows for safety as the driver can release the page when it is safe to do so. The vFree strategy has been shown to significantly reduce the performance overhead of the vIOMMU, particularly in network-heavy workloads. However, most of the existing research on vFree has focused on single VM workloads, and there is a lack of comprehensive evaluation of these strategies in multi-VM environments. This paper aims to fill this gap by evaluating the performance of vFree in a multi-VM environment and comparing it to traditional approaches for managing the vIOMMU. By doing so, we can gain insights into the effectiveness of vFree in real-world scenarios where multiple VMs are running concurrently and sharing resources.
